\chapter{Kiến trúc hệ thống}
\label{chap:architecture}

Chương này trình bày kiến trúc chi tiết của hệ thống tự chẩn đoán sự cố mạng truyền thông dựa trên hướng tiếp cận Self-Evolving AI Agents tích hợp trên nền tảng NIKA. Kiến trúc hệ thống được thiết kế hướng tới khả năng mô phỏng môi trường mạng thực tế, hỗ trợ nhiều luồng suy luận của agent, và đặc biệt là hai mô-đun tự tiến hóa: Mô-đun bộ nhớ quy trình (Memory Module) và Mô-đun tiến hóa công cụ (Tool Evolution Module).

\section{Kiến trúc tổng quan}
\label{sec:overall_arch}

Kiến trúc tổng quan của hệ thống bao gồm 4 thành phần chính được kết nối chặt chẽ với nhau như được mô tả trong Hình dưới đây (triển khai dạng luồng tương tác):
\begin{enumerate}
    \item **Môi trường mô phỏng mạng (Network Environment):** Chạy trên nền tảng Kathará ảo hóa dưới dạng các Docker container đại diện cho thiết bị mạng (routers, switches, hosts) kết nối theo các topology phức tạp.
    \item **AI Agent và Workflow chẩn đoán:** Đóng vai trò là thực thể thực hiện chẩn đoán sự cố, hỗ trợ các workflow suy luận phổ biến như ReAct, Plan \& Execute, và Reflexion.
    \item **Hệ thống MCP Servers:** Cung cấp cầu nối tiêu chuẩn hóa để AI Agent có thể truy vấn thông tin, thực thi lệnh chẩn đoán trên thiết bị mạng ảo hóa, và gửi kết quả chẩn đoán.
    \item **Mô-đun Tự Tiến Hóa (Self-Evolution Modules):** Gồm mô-đun bộ nhớ quy trình (lưu trữ kinh nghiệm chẩn đoán lâu dài) và mô-đun tiến hóa công cụ (tự sinh công cụ chẩn đoán mới).
\end{enumerate}

\section{Môi trường mạng ảo hóa và các dịch vụ MCP}
\label{sec:net_env_mcp}

Để AI Agent có khả năng tương tác linh hoạt với môi trường mạng giả lập mà không phụ thuộc vào một cấu trúc phần cứng cụ thể, hệ thống ứng dụng giao thức Model Context Protocol (MCP) do Anthropic đề xuất. Nền tảng NIKA thiết kế và khởi chạy 6 máy chủ MCP chuyên biệt trong thư mục \texttt{src/nika/service/mcp\_server}:
\begin{itemize}
    \item **Kathará Base MCP Server:** Cung cấp các công cụ chẩn đoán cơ bản như kiểm tra kết nối giữa các host (\texttt{get\_reachability}), ping một cặp thiết bị (\texttt{ping\_pair}), đo băng thông (\texttt{iperf\_test}), kiểm tra dịch vụ (\texttt{systemctl\_ops}), xem cấu hình giao tiếp mạng (\texttt{get\_host\_net\_config}), thống kê giao thức (\texttt{netstat}) và thực thi shell command trong container (\texttt{exec\_shell}).
    \item **FRR MCP Server:** Tương tác với phần mềm định tuyến FRRouting chạy trên các router để trích xuất cấu hình định tuyến động OSPF (\texttt{frr\_get\_ospf\_conf}), BGP (\texttt{frr\_get\_bgp\_conf}) và bảng định tuyến hiện tại (\texttt{frr\_show\_ip\_route}).
    \item **BMv2 MCP Server:** Tương tác với các P4 software switches chạy phần mềm BMv2 để đọc dữ liệu telemetry, đọc giá trị counter (\texttt{bmv2\_counter\_read}), dump các bảng luồng hành vi (\texttt{bmv2\_table\_dump}) và đọc thanh ghi (\texttt{bmv2\_register\_read}).
    \item **Telemetry MCP Server:** Truy vấn cơ sở dữ liệu InfluxDB thu thập từ In-band Network Telemetry (INT) để lấy các chỉ số đo lường hiệu năng mạng theo thời gian thực.
    \item **Task Management MCP Server:** Cung cấp công cụ \texttt{submit} để agent gửi kết quả chẩn đoán cuối cùng (gồm trạng thái lỗi, thiết bị bị lỗi và phân tích nguyên nhân gốc rễ RCA).
    \item **Evolved Diagnostic Toolbox MCP Server:** Cung cấp công cụ khám phá (\texttt{list\_evolved\_tools}) và thực thi (\texttt{execute\_evolved\_tool}) các công cụ chẩn đoán mới được tiến hóa thành công trong thư viện.
\end{itemize}

\section{Các luồng suy luận chẩn đoán của AI Agent}
\label{sec:agent_workflows}

AI Agent trong hệ thống đóng vai trò nhận mô tả tác vụ từ môi trường và sử dụng các công cụ chẩn đoán MCP để tìm ra thiết bị lỗi và nguyên nhân gốc rễ. Hệ thống hỗ trợ ba cơ chế workflow chẩn đoán chính:
\begin{itemize}
    \item **ReAct (Reasoning and Acting):** Agent luân phiên thực hiện bước suy luận ngữ cảnh (Thought) và gọi công cụ chẩn đoán tương ứng (Action) cho đến khi thu thập đủ bằng chứng chẩn đoán để đưa ra quyết định cuối cùng.
    \item **Plan \& Execute:** Một bộ lập kế hoạch (Planner) sẽ phân tích bài toán và sinh ra một chuỗi các bước chẩn đoán cần thiết. Một bộ thực thi (Executor) sẽ chạy tuần tự các bước đó. Nếu có lỗi phát sinh hoặc kết quả bất thường, bộ tái lập kế hoạch (Replanner) sẽ điều chỉnh danh sách công việc.
    \item **Reflexion:** Tích hợp vòng lặp phản tư. Sau mỗi lần chẩn đoán thất bại, agent tự phân tích vết suy luận lỗi, rút kinh nghiệm dưới dạng episodic memory và thực hiện chẩn đoán lại với định hướng chiến lược mới nhằm tránh lặp lại sai lầm.
\end{itemize}

\section{Mô-đun Bộ Nhớ Quy Trình Tiến Hóa (Memory Module)}
\label{sec:memory_module}

Mô-đun Bộ nhớ quy trình giúp lưu trữ tri thức lâu dài cho agent qua các incident chẩn đoán liên tiếp. Mô-đun này được kế thừa và kết hợp các ưu điểm từ 3 công trình lớn: **MemInsight**, **LightMem**, và **A-Mem**.

\subsection{Cơ chế lưu trữ hỗn hợp}
Bộ nhớ quy trình được lưu trữ dưới hai dạng:
\begin{enumerate}
    \item **PostgreSQL (Canonical Store):** Lưu trữ chính thức các bản ghi bộ nhớ, thông tin nguồn gốc (provenance), mức độ tin cậy (confidence), phiên bản, và các liên kết đồ thị giữa các ghi chú.
    \item **Qdrant (Semantic Index):** Lưu trữ vector nhúng của nội dung ghi nhớ để hỗ trợ tìm kiếm tương đồng ngữ nghĩa. Khi cơ sở dữ liệu Qdrant chưa sẵn sàng, hệ thống tự động fallback sang tìm kiếm văn bản (lexical/full-text search) trên PostgreSQL.
\end{enumerate}

\subsection{Truy xuất nhận biết thuộc tính (MemInsight-style Retrieval)}
Trước mỗi tập chẩn đoán (episode), từ ngữ cảnh công khai ban đầu (public task context), hệ thống sinh ra một tập thuộc tính truy vấn (\texttt{query\_attributes}) gồm: tên kịch bản mạng, lớp topology, giao thức sử dụng (OSPF, BGP, DHCP, etc.), dịch vụ kích hoạt.
Bộ nhớ được truy xuất theo công thức tính điểm kết hợp:
\begin{equation}
\text{Score} = \text{Similarity}_{\text{semantic}} + \text{Bonus}_{\text{attribute}} + \text{Weight}_{\text{confidence}} - \text{Penalty}_{\text{redundancy}}
\end{equation}
Điều này giúp agent lọc bỏ các kinh nghiệm chẩn đoán không liên quan đến cấu trúc mạng hiện tại (ví dụ: không lấy kinh nghiệm BGP Clos khi đang chẩn đoán OSPF Enterprise).

\subsection{Nén vết tương tác (LightMem-style Compaction)}
Sau khi tập chẩn đoán kết thúc, thay vì lưu trữ toàn bộ lịch sử hội thoại dài và nhiều nhiễu (\texttt{messages.jsonl}), mô-đun thực hiện quy trình nén vết chẩn đoán:
\begin{itemize}
    \item Loại bỏ các thông tin nhạy cảm hoặc định danh cụ thể (redact IPv4/IPv6 thành \texttt{<ip>}, MAC thành \texttt{<mac>}, tên thiết bị thành \texttt{<device\_A>}).
    \item Phân đoạn lịch sử theo từng pha chẩn đoán (ví dụ: kiểm tra ban đầu, kiểm tra định tuyến, xác định lỗi).
    \item LLM đóng vai trò nén thông tin, trích xuất thành các ghi chú nguyên tử ngắn gọn (atomic notes).
\end{itemize}

\subsection{Đồ thị liên kết và Cơ chế lọc điểm số (A-Mem \& Score-gated Validation)}
Các ghi chú nguyên tử mới sinh ra được so sánh với bộ nhớ hiện tại để tạo các liên kết đồ thị:
\begin{itemize}
    \item **supports:** Củng cố ghi chú cũ.
    \item **refines:** Chi tiết hóa hoặc bổ sung điều kiện ràng buộc cho ghi chú cũ.
    \item **contradicts:** Mâu thuẫn với ghi chú cũ (xảy ra khi môi trường hoặc cấu hình mạng thay đổi).
\end{itemize}

Để đảm bảo an toàn và tính khách quan cho benchmark (tránh rò rỉ đáp án), việc cập nhật bộ nhớ phải tuân theo cơ chế **Score-gated Validation**:
\begin{itemize}
    \item Chỉ khi NIKA Evaluator chấm điểm chẩn đoán thành công tuyệt đối (độ chính xác RCA = 1.0), các ghi chú quy trình chẩn đoán (procedural notes) mới được gán nhãn \texttt{validated} và tăng độ tin cậy.
    \item Nếu chẩn đoán thất bại hoặc thành công một phần, hệ thống chỉ chấp nhận lưu trữ các ghi chú dạng cảnh báo (avoid rules hoặc caution notes, ví dụ: "Không được kết luận lỗi DNS chỉ dựa vào kiểm tra curl"), đồng thời giảm độ tin cậy của các ghi chú đã hướng dẫn sai.
\end{itemize}

\section{Mô-đun Tiến Hóa Công Cụ (Tool Evolution Module)}
\label{sec:tool_evolution}

Mô-đun Tiến hóa công cụ cho phép AI Agent tự mở rộng thư viện công cụ của mình để vượt qua các giới hạn năng lực tại test-time. Thư viện công cụ lưu trữ dưới dạng \texttt{state.json} gồm 3 thành phần chính:

\subsection{Mastery mô tả công cụ (DRAFT-like Refinement)}
Đối với các công cụ nguyên thủy (primitive tools), hệ thống không thay đổi mã nguồn thực thi mà cập nhật một lớp phủ tài liệu hướng dẫn (\texttt{ToolMastery}) dựa trên kết quả chạy thực tế:
\begin{itemize}
    \item **preconditions:** Điều kiện phù hợp nhất để gọi công cụ.
    \item **parameter\_guidance:** Hướng dẫn truyền tham số đúng kiểu và đúng ngữ cảnh thiết bị.
    \item **failure\_semantics:** Ý nghĩa của các mã lỗi trả về.
\end{itemize}

\subsection{Tổng hợp quy trình phức hợp và sinh mã nguồn Python (TTE \& Alita-style)}
Khi phát hiện một thiếu hụt năng lực chẩn đoán (capability gap), agent có thể đề xuất 2 dạng công cụ mới:
\begin{enumerate}
    \item **Composite Workflow:** Chuỗi các bước gọi tuần tự các công cụ nguyên thủy an toàn, được tham số hóa để tái sử dụng (ví dụ: kiểm tra định tuyến tuần tự từ Host -> Router -> Gateway).
    \item **Generated Python Helper:** Hàm Python thuần thực hiện các tính toán phức tạp trên dữ liệu chẩn đoán thu thập được (ví dụ: tính toán tỷ lệ mất gói từ các số liệu đếm counter).
\end{enumerate}

\subsection{Quy trình xác thực an toàn (AST \& Sandbox Validation)}
Để tránh rủi ro bảo mật hệ thống và đảm bảo công cụ sinh ra hoạt động chính xác, toàn bộ các công cụ Python mới sinh ra đều phải đi qua pipeline xác thực nghiêm ngặt:
\begin{enumerate}
    \item **Xác thực cấu trúc (AST Verification):** Sử dụng thư viện \texttt{ast} của Python để kiểm tra cấu trúc cú pháp, đảm bảo không sử dụng các hàm cấm (như \texttt{eval}, \texttt{exec}, \texttt{open}, \texttt{\_\_import\_\_}) và chỉ import các thư viện toán học trong allowlist.
    \item **Thực thi Sandbox (Docker Sandbox execution):** Chạy thử nghiệm công cụ trong một container Docker cô lập hoàn toàn mạng (\texttt{--network none}) với các tham số mẫu. Hàm sinh ra phải chạy thành công và trả về định dạng JSON hợp lệ.
    \item **Xác thực kết quả thực nghiệm:** Công cụ mới sinh ra ban đầu ở trạng thái thử nghiệm (\texttt{candidate}). Nó chỉ được thăng cấp lên chính thức (\texttt{promoted}) sau khi hỗ trợ agent chẩn đoán thành công sự cố ở ít nhất 2 ngữ cảnh kịch bản mạng khác nhau.
\end{enumerate}